\documentclass[aspectratio=169,10pt]{beamer}
\usetheme[numbering=fraction,progressbar=frametitle]{metropolis}
\usepackage{fontspec}
\setsansfont{Latin Modern Sans}
\usepackage{amsmath,amssymb,mathtools,graphicx,booktabs}
\usepackage{polyglossia}
\setdefaultlanguage{french}
\setbeamertemplate{navigation symbols}{}
\setbeamerfont{frametitle}{size=\large}
\setbeamerfont{normal text}{size=\normalsize}
\setlength{\parskip}{.45em}
\hypersetup{pdfauthor={},pdfsubject={Lecture slides},pdfkeywords={},pdftitle={Statistiques descriptives à une variable}}
\providecommand{\N}{\mathbb{N}}
\providecommand{\Z}{\mathbb{Z}}
\providecommand{\D}{\mathbb{D}}
\providecommand{\Q}{\mathbb{Q}}
\providecommand{\R}{\mathbb{R}}
\providecommand{\C}{\mathbb{C}}
\providecommand{\vect}[1]{\overrightarrow{#1}}
\providecommand{\vi}{\vec{\imath}}
\providecommand{\vj}{\vec{\jmath}}
\providecommand{\vk}{\vec{k}}
\providecommand{\norm}[1]{\left\lVert #1\right\rVert}
\providecommand{\vc}[2]{\begin{pmatrix}#1\\#2\end{pmatrix}}
\providecommand{\deter}[2]{\det\left(#1,#2\right)}
\providecommand{\abs}[1]{\left\lvert #1\right\rvert}
\providecommand{\dd}{\,\mathrm{d}}
\providecommand{\e}{\mathrm{e}}
\providecommand{\Card}{\operatorname{Card}}
\providecommand{\E}{\operatorname{E}}
\providecommand{\V}{\operatorname{V}}
\providecommand{\degre}{^{\circ}}
\providecommand{\un}[1]{\,\mathrm{#1}}
\providecommand{\sci}[2]{#1\times10^{#2}}
\providecommand{\pgcd}{\operatorname{PGCD}}
\providecommand{\ppcm}{\operatorname{PPCM}}
\title{Statistiques descriptives à une variable}
\subtitle{Cours et exemples}
\date{}
\author{}
\begin{document}
\begin{frame}[plain]
\vfill
{\usebeamerfont{title}\inserttitle\par}
\medskip
{\usebeamerfont{subtitle}\insertsubtitle\par}
\medskip
{\usebeamercolor[fg]{progress bar}\rule{\textwidth}{.5pt}}
\vfill
\end{frame}
\begin{frame}{Effectifs, fréquences et moyenne}
Une série décrit un caractère sur une population. Si la valeur \(x_i\) apparaît \(n_i\) fois, \(n_i\) est son effectif. L'effectif total non nul est \(N=\sum_{i=1}^{p}n_i\) ; sa fréquence est \(f_i=n_i/N\). La moyenne est
\[\overline x=\dfrac{\sum_{i=1}^{p}n_ix_i}{N}=\sum_{i=1}^{p}f_ix_i.\]
L'étendue est \(x_{\max}-x_{\min}\).
\end{frame}
\begin{frame}{Effectifs, fréquences et moyenne}
On étudie les pointures de \(30\) élèves :
  \begin{center}\begin{tabular}{lllllllllllllll}
36 & 42 & 40 & 41 & 41 & 38 & 39 & 43 & 42 & 42 & 42 & 40 & 36 & 40 & 39 \\ 
41 & 40 & 38 & 44 & 42 & 40 & 42 & 43 & 38 & 43 & 36 & 41 & 41 & 44 & 43 \\ 
\end{tabular}\end{center} 
Compléter les effectifs, calculer les fréquences, la moyenne et l'étendue.
\end{frame}
\begin{frame}{Correction : Effectifs, fréquences et moyenne}
Effectifs : \(3,0,3,2,5,5,6,4,2\), de somme \(30\). Les fréquences sont ces effectifs divisés par \(30\), soit en pourcentage \(10\) ; \(0\) ; \(10\) ; environ \(6{,}67\) ; \(16{,}67\) ; \(16{,}67\) ; \(20\) ; \(13{,}33\) ; \(6{,}67\).
\(\sum n_ix_i=1217\), donc \(\overline x=1217/30\approx40{,}57\). L'étendue vaut \(8\). La moyenne n'est pas forcément une valeur de la série. La somme des pourcentages arrondis peut différer légèrement de \(100\).
\end{frame}
\begin{frame}{Médiane et effectifs cumulés}
Une médiane est un nombre tel qu'au moins la moitié des valeurs lui soient inférieures ou égales, et au moins la moitié supérieures ou égales.
Dans la série ordonnée de \(N\) valeurs, si \(N\) est impair, prendre le rang \((N+1)/2\) ; si \(N\) est pair, choisir la moyenne des valeurs de rangs \(N/2\) et \(N/2+1\).
\end{frame}
\begin{frame}{Médiane et effectifs cumulés}
Avec des répétitions, plus de \(50\,\%\) des valeurs peuvent être inférieures ou égales à la médiane.
\end{frame}
\begin{frame}{Médiane et effectifs cumulés}
Compléter les effectifs cumulés croissants (ECC) des pointures et déterminer la médiane.
\end{frame}
\begin{frame}{Correction : Médiane et effectifs cumulés}
Les rangs centraux sont \(15\) et \(16\). Les ECC passent de \(13\) pour \(40\) à \(18\) pour \(41\) : les deux valeurs centrales valent \(41\), donc \(M_e=41\). Au moins la moitié des pointures sont inférieures ou égales à \(41\), et au moins la moitié supérieures ou égales à \(41\).
\end{frame}
\begin{frame}{Médiane et effectifs cumulés}
Calculer les médianes de \(1;2;2;8;9\) et de \(1;2;8;9\).
\end{frame}
\begin{frame}{Correction : Médiane et effectifs cumulés}
La troisième valeur de la première série vaut \(2\). Dans la seconde, \((2+8)/2=5\) ; la médiane choisie n'est pas une valeur de la série.
\end{frame}
\begin{frame}{Quartiles et diagramme en boîte}
Le premier quartile \(Q_1\) est la valeur de rang \(\lceil N/4\rceil\) : la plus petite valeur de la série telle qu'au moins \(25\,\%\) des valeurs lui soient inférieures ou égales. \(Q_3\) est la valeur de rang \(\lceil3N/4\rceil\).
L'écart interquartile est \(Q_3-Q_1\). Le diagramme en boîte indique minimum, \(Q_1\), médiane, \(Q_3\), maximum.
\end{frame}
\begin{frame}{Quartiles et diagramme en boîte}
\begin{columns}[T,onlytextwidth]
\begin{column}{.57\textwidth}
Déterminer \(Q_1\), \(Q_3\) et l'écart interquartile des pointures, puis tracer le diagramme en boîte.
\end{column}
\begin{column}{.4\textwidth}
\centering
\includegraphics[width=\linewidth,height=4.5cm,keepaspectratio]{figures/seconde-07-1.pdf}
\end{column}
\end{columns}
\end{frame}
\begin{frame}{Correction : Quartiles et diagramme en boîte}
\(N/4=7{,}5\) : rang \(8\), donc \(Q_1=39\). \(3N/4=22{,}5\) : rang \(23\), donc \(Q_3=42\). Écart interquartile \(3\). Le diagramme utilise \((36;39;41;42;44)\).
\end{frame}
\begin{frame}{Quartiles et diagramme en boîte}
Déterminer les quartiles de \(1;2;2;3;4;4;4;5\). Y a-t-il exactement \(25\,\%\) des valeurs inférieures ou égales à \(Q_1\) ?
\end{frame}
\begin{frame}{Correction : Quartiles et diagramme en boîte}
\(Q_1=2\) (rang \(2\)), \(Q_3=4\) (rang \(6\)). Ici \(3/8=37{,}5\,\%\) des valeurs sont inférieures ou égales à \(Q_1\). Il faut dire « au moins ».
\end{frame}
\begin{frame}{Ajout ou suppression d'une valeur}
Ajouter \(a\) à une série de taille \(N\) donne \(\overline x'=(N\overline x+a)/(N+1)\).
Retirer une occurrence de \(a\), pour \(N>1\), donne \(\overline x'=(N\overline x-a)/(N-1)\). Utiliser la somme exacte, sans arrondi intermédiaire. Réexaminer la médiane dans la nouvelle série ordonnée.
\end{frame}
\begin{frame}{Ajout ou suppression d'une valeur}
Dans la série des pointures, la somme vaut \(1217\), la médiane \(41\).

\begin{enumerate}
\setcounter{enumi}{0}\item Ajouter une pointure \(36\) et calculer les nouveaux indicateurs.
\item Repartir de la série initiale, retirer une pointure \(44\) et recommencer.

\end{enumerate}
\end{frame}
\begin{frame}{Correction : Ajout ou suppression d'une valeur}
1. Somme \(1253\), effectif \(31\) : \(\overline x'=1253/31\approx40{,}42\). Le rang central \(16\) donne toujours \(41\). 2. Somme \(1173\), effectif \(29\) : \(\overline x''=1173/29\approx40{,}45\). Le rang central \(15\) donne \(41\). La moyenne diminue dans les deux cas ; la médiane reste inchangée. Les situations sont indépendantes.
\end{frame}
\begin{frame}{Ajout ou suppression d'une valeur}
Justifier qu'ajouter la moyenne ne change pas la moyenne. L'étendue peut-elle diminuer lors d'un ajout ?
\end{frame}
\begin{frame}{Correction : Ajout ou suppression d'une valeur}
Si \(a=\overline x\), \((N\overline x+\overline x)/(N+1)=\overline x\). Lors d'un ajout, le minimum ne peut qu'être conservé ou diminuer, le maximum être conservé ou augmenter : l'étendue ne diminue pas. Lors d'une suppression, elle n'augmente pas.
\end{frame}
\begin{frame}{Variance et écart type}
\(V=\dfrac{\sum n_i(x_i-\overline x)^2}{N}\) est la variance ; \(\sigma=\sqrt V\) est l'écart type. Il mesure la dispersion autour de la moyenne et a la même unité que les données. On peut aussi calculer \(V=\dfrac{\sum n_ix_i^2}{N}-\overline x^{\,2}\).
\end{frame}
\begin{frame}{Variance et écart type}
Calculer la variance et l'écart type des pointures.
\end{frame}
\begin{frame}{Correction : Variance et écart type}
\(\sum n_ix_i^2=49519\), donc
\[V=\dfrac{49519}{30}-\left(\dfrac{1217}{30}\right)^2=\dfrac{4481}{900}\approx4{,}979.\]
\(\sigma=\sqrt{4481}/30\approx2{,}23\). L'écart type n'est pas la distance maximale à la moyenne.
\end{frame}
\begin{frame}{Variance et écart type}
Comparer \(A:8;10;12\) et \(B:0;10;20\) avec les couples moyenne-écart type et médiane-écart interquartile.
\end{frame}
\begin{frame}{Correction : Variance et écart type}
Les moyennes et médianes valent toutes \(10\). \(V_A=8/3\), \(\sigma_A\approx1{,}63\) ; \(V_B=200/3\), \(\sigma_B\approx8{,}16\). Quartiles : \(8,12\) pour \(A\) ; \(0,20\) pour \(B\). Écarts interquartiles \(4\) et \(20\). \(B\) est beaucoup plus dispersée malgré des centres identiques.
\end{frame}
\begin{frame}{Linéarité de la moyenne}
Remplacer chaque \(x_i\) par \(kx_i+c\) donne une moyenne \(k\overline x+c\), un écart type \(|k|\sigma\) et une variance \(k^2V\). Ajouter une constante ne change ni variance ni écart type.
\end{frame}
\begin{frame}{Linéarité de la moyenne}
Moyenne initiale \(12/20\), écart type \(3\). Calculer les nouveaux indicateurs pour trois modifications indépendantes, sans plafonnement des notes : ajouter \(1{,}5\) ; multiplier par \(0{,}8\) ; multiplier par \(1{,}1\) puis ajouter \(0{,}5\).
\end{frame}
\begin{frame}{Correction : Linéarité de la moyenne}
Moyennes : \(13{,}5\), \(9{,}6\), \(13{,}7\). Écarts types : \(3\), \(2{,}4\), \(3{,}3\). Multiplier les notes par \(1{,}1\) les augmente ; cela ne signifie pas augmenter la difficulté de l'évaluation. Avec un plafonnement à \(20\), la transformation ne serait plus affine sur toutes les données.
\end{frame}
\begin{frame}[fragile]{Python}
\small
\begin{verbatim}
from math import sqrt

def moyenne_ecart_type(valeurs):
    n = len(valeurs)
    if n == 0:
        raise ValueError("liste vide")
    m = sum(valeurs)/n
    v = sum((x-m)**2 for x in valeurs)/n
    return m, sqrt(v)
\end{verbatim}
\end{frame}
\begin{frame}{Linéarité de la moyenne}
Expliquer le programme ; prévoir son résultat pour \texttt{[8, 10, 12]}.
\end{frame}
\begin{frame}{Correction : Linéarité de la moyenne}
Il vérifie que la liste est non vide, calcule la moyenne, la moyenne des carrés des écarts, puis sa racine. Résultat : \((10;\sqrt{8/3})\approx(10;1{,}633)\).
\end{frame}
\begin{frame}{Séries continues regroupées en classes}
Une variable quantitative continue peut prendre toutes les valeurs d'un intervalle. Ici, les observations sont regroupées en classes de même amplitude ; le centre de \([a;b[\) vaut \((a+b)/2\).
Si les moyennes \(m_i\) des classes sont connues, la moyenne globale exacte est \(\sum n_im_i/N\). À défaut, les remplacer par les centres donne une estimation justifiée par un modèle uniforme dans chaque classe.
\end{frame}
\begin{frame}{Séries continues regroupées en classes}
Compléter, estimer la moyenne à partir des centres et déterminer la classe médiane.
\end{frame}
\begin{frame}{Correction : Séries continues regroupées en classes}
Amplitude \(10\) cm, effectif total \(100\). Fréquences en pourcentage : \(8,22,35,28,7\). Centres : \(155,165,175,185,195\) cm. La somme pondérée des centres vaut \(17540\), donc la moyenne estimée est \(175{,}4\) cm. Les rangs centraux \(50,51\) se trouvent dans \([170;180[\), classe médiane.
\end{frame}
\begin{frame}{Séries continues regroupées en classes}
Si les moyennes exactes des classes étaient \(154;166;174;184;196\) cm, quelle serait la moyenne globale ?
\end{frame}
\begin{frame}{Correction : Séries continues regroupées en classes}
\((8\times154+22\times166+35\times174+28\times184+7\times196)/100=174{,}98\) cm.
\end{frame}
\begin{frame}{Histogramme et fréquences cumulées}
L'aire de chaque rectangle d'un histogramme est proportionnelle à l'effectif. À amplitudes égales, les hauteurs peuvent être les effectifs. Le polygone des fréquences cumulées joint les points (borne supérieure ; fréquence cumulée), en commençant à la première borne inférieure avec \(0\).
\end{frame}
\begin{frame}{Histogramme et fréquences cumulées}
\begin{columns}[T,onlytextwidth]
\begin{column}{.57\textwidth}
Construire les deux graphiques pour les tailles de l'exemple précédent.
\end{column}
\begin{column}{.4\textwidth}
\centering
\includegraphics[width=\linewidth,height=2.25cm,keepaspectratio]{figures/seconde-07-2.pdf}
\includegraphics[width=\linewidth,height=2.25cm,keepaspectratio]{figures/seconde-07-3.pdf}
\end{column}
\end{columns}
\end{frame}
\begin{frame}{Histogramme et fréquences cumulées}
Sous l'hypothèse d'uniformité dans \([170;180[\), la fréquence cumulée passe linéairement de \(30\,\%\) à \(65\,\%\). Rechercher l'abscisse correspondant à \(50\,\%\).
\end{frame}
\begin{frame}{Histogramme et fréquences cumulées}
Calculer cette estimation et préciser sa portée.
\end{frame}
\begin{frame}{Correction : Histogramme et fréquences cumulées}
On résout \((m-170)/10=(50-30)/(65-30)\), d'où \(m=170+10\times20/35\approx175{,}71\) cm. Il s'agit d'une estimation sous hypothèse d'uniformité, pas de la médiane exacte des données individuelles. Les effectifs des classes seuls ne déterminent pas cette dernière.
\end{frame}
\end{document}
